% Imporditud: tools/import_eksperimendid.py
% Allikas: Eesti füüsikaolümpiaadi eksperimendi ülesannete
%   kogu (Rael Kalda, 2023), ülesanne Ü122, lahendus L122.
\setAuthor{EFO žürii}
\setRound{lõppvoor}
\setYear{2021}
\setNumber{P E1}
\setDifficulty{3}
\setTopic{geomeetriline optika}
\setVahendid{vesi, toiduõli, läbipaistev pudel ära lõigatud ülemise osaga}
\setPunktid{10}
\setAllikas{Eksperimendi kogu Ü122, lahendus lk 97}
\setLahendusseis{toores}

\prob{Optiline tihedus}
Hinnake, kas optiline tihedus on suurem veel või toiduõlil. Põhjendage tulemust.

\hint

\solu
Tuleb vaadata läbi vedelike mingit eset või kujundit paberil. Asetades need pudelile piisavalt lähedale, näeme suurendatud (ümber pööramata) kujutise. Mida suurem on kujutis, seda suurem on optiline tihedus, sest murdva pinna kumerus on mõlemal juhul ühesugune. Suuremale murdumisnäitajale vastab suurem murdumisnurk, sellele aga suurem kujutis (parim seletus on joonise abil). Kuivõid vesi ja toiduõli omavahel ei segune, võime need paremaks vaatlemiseks korraga pudelisse valada (vesi jääb alla ja toiduõli ülesse). Näeme, et toiduõlis on kujutis suurem, mistõttu on toiduõli optiliselt tihedam. Vee murdumisnäitaja on 1,33, toiduõli oma 1,47.
\probend
